% !TEX root = ../main.tex
\chapter{Simulación}
\label{ch:simulacion}

El problema plano del capítulo~\ref{ch:modelos-previos} permite comprobar
el cálculo numérico antes de abordar Schwarzschild. Conocemos su solución
exacta: podemos calcular el campo mediante un programa y medir cuánto se
aparta de esa referencia. En este capítulo presentamos el procedimiento,
los experimentos ejecutados y lo que permiten concluir.

\section{Problema y pruebas realizadas}

Utilizamos las variables adimensionales definidas en
\eqref{eq:kg-adimensional}. Los sombreros recuerdan que hemos dividido las
magnitudes físicas por sus escalas de referencia. La ecuación que resuelve
el programa es
\[
\widehat\Phi_{\widehat t\widehat t}
=\widehat\Phi_{\widehat x\widehat x}-\widehat\mu^2\widehat\Phi.
\]
Trabajamos en $0\le\widehat x<2\pi$, con $\widehat\mu=0.7$, amplitud
inicial uno y tiempo final $\widehat t=20$. Para cada entero $j$ elegido,
el perfil inicial y su velocidad son
\[
\widehat\Phi(\widehat x,0)=\cos(j\widehat x),\qquad
\widehat\Phi_{\widehat t}(\widehat x,0)=0.
\]
El índice $j$ cuenta cuántas oscilaciones completas del coseno caben en el
dominio. Se mantiene el contorno periódico del modelo previo: el extremo
derecho se conecta con el izquierdo. Esta elección hace compatible el
perfil con el dominio y permite utilizar la solución exacta; no representa
las fronteras de un agujero negro.

Se han realizado dos pruebas. En la primera se usan $256$ puntos y se
varía $j=1,2,3,4$ para comprobar la relación entre número de onda y
frecuencia. En la segunda se fija $j=2$ y se repite el cálculo con $64$,
$128$, $256$ y $512$ puntos para comprobar cómo disminuye el error.
El caso $j=2$ con $256$ puntos es común a ambas pruebas: en total son
siete evoluciones distintas.

Los parámetros son elecciones de validación, sin una masa física en SI
seleccionada. Los datos, tablas y figuras de las siete ejecuciones se
conservan en \texttt{results/flat/}; el apartado~\ref{sec:reproduccion-plana}
explica cómo reproducirlos.

% !TEX root = ../main.tex
\section{Cómo se calcula la evolución}
\label{sec:metodo-plano}

\subsection{Aproximar el campo y sus derivadas}
El programa almacena el campo en $N$ posiciones igualmente separadas.
La distancia entre dos posiciones consecutivas es $h=\widehat L/N$.
Se almacena el extremo izquierdo y se omite la copia del extremo derecho,
porque ambos representan el mismo punto del dominio periódico. Para
calcular junto a un extremo se utiliza como vecino el punto correspondiente
del extremo opuesto.

La ecuación necesita una segunda derivada espacial, pero el programa
dispone de valores en posiciones separadas. La aproximamos utilizando el
valor de un punto y los de sus dos vecinos. Su origen se ve al desarrollar
por Taylor una función suave $f$:
\[
\begin{aligned}
f(x+h)&=f(x)+hf'(x)+\tfrac12h^2f''(x)+\cdots,\\
f(x-h)&=f(x)-hf'(x)+\tfrac12h^2f''(x)+\cdots.
\end{aligned}
\]
Al sumar, los términos de la primera derivada se cancelan. Despejar la
segunda derivada conduce a
\[
f''(x)\simeq\frac{f(x+h)-2f(x)+f(x-h)}{h^2}.
\]
Esta fórmula compara el valor central con los de alrededor y aproxima
cómo cambia la pendiente del perfil. Se llama diferencia centrada porque
utiliza un vecino a cada lado. Para una función suficientemente suave,
el primer error que queda después de la cancelación es proporcional a
$h^2$ al hacer $h$ pequeño. En los perfiles sinusoidales utilizados se
cumple esa regularidad.

\subsection{Avanzar de un instante al siguiente}
El tiempo también se divide en pasos de tamaño $\Delta\widehat t$.
Escribimos $\Phi_i^n$ para el valor numérico del campo adimensional en
la posición $i$ y en el instante $n\Delta\widehat t$. El subíndice $i$
identifica una posición; el superíndice $n$ identifica un instante, no
una potencia.

Aproximamos la segunda derivada temporal del mismo modo, utilizando el
instante anterior, el actual y el siguiente. Al sustituir ambas
aproximaciones en Klein--Gordon y despejar el valor siguiente resulta
\begin{equation}
\begin{aligned}
\Phi_i^{n+1}={}&2\Phi_i^n-\Phi_i^{n-1}\\
&+(\Delta\widehat t)^2
\left(\frac{\Phi_{i+1}^n-2\Phi_i^n+\Phi_{i-1}^n}{h^2}
-\widehat\mu^2\Phi_i^n\right).
\end{aligned}
\label{eq:esquema-plano}
\end{equation}
La expresión entre paréntesis aproxima la segunda derivada temporal que
predice la EDP. Así, cada paso utiliza el campo actual, sus vecinos y el
campo del instante anterior. La regla es explícita: todos esos valores
ya se conocen cuando se calcula el siguiente.

Para empezar necesitamos dos instantes. El primero se fija con el perfil
inicial. El segundo se aproxima mediante Taylor en el tiempo. Como la
velocidad inicial es cero, el término lineal desaparece y queda
\begin{equation}
\Phi_i^1=\Phi_i^0+\frac{(\Delta\widehat t)^2}{2}
\left(\frac{\Phi_{i+1}^0-2\Phi_i^0+\Phi_{i-1}^0}{h^2}
-\widehat\mu^2\Phi_i^0\right).
\label{eq:inicio-plano}
\end{equation}
Esto incorpora el cambio inicial que exige la ecuación aunque el campo
parta con velocidad nula. A partir de esos dos instantes se repite
\eqref{eq:esquema-plano} hasta el tiempo final.

\subsection{Elección del paso y comprobación del error}
El paso temporal debe elegirse junto con la separación espacial. Un paso
demasiado grande puede hacer crecer errores del cálculo aunque la
solución exacta permanezca acotada. Para este esquema periódico utilizamos
la restricción de estabilidad
\begin{equation}
\Delta\widehat t<\frac{2}{\sqrt{\widehat\mu^2+4/h^2}}.
\label{eq:estabilidad-plana}
\end{equation}
Es una condición del procedimiento numérico. En las ejecuciones se propone
un paso igual a $0.8$ veces esa cota y se reduce ligeramente, si hace falta,
para que un número entero de pasos termine exactamente en
$\widehat t=20$. Al aumentar $N$ también disminuye el paso temporal.

Las diferencias centradas aproximan ambas derivadas con errores de orden
dos para soluciones suaves. Por ello esperamos que el error disminuya
aproximadamente por un factor cuatro al reducir a la mitad los pasos
espacial y temporal. Esta expectativa se contrasta con las cuatro mallas
utilizadas; los resultados y el alcance de esa comprobación se explican
a continuación.

% !TEX root = ../main.tex
% Resultados de las siete ejecuciones identificadas en results/flat/manifest.json.
\section{Qué muestran los resultados}
\label{sec:validacion-plana}

\subsection{Comparación con la solución exacta}
Para los datos elegidos, la referencia del capítulo~\ref{ch:modelos-previos}
se escribe
\[
\widehat\Phi_{\rm ex}(\widehat x,\widehat t)
=\cos(j\widehat x)\cos(\widehat\omega_j\widehat t),\qquad
\widehat\omega_j=\sqrt{j^2+\widehat\mu^2}.
\]
El programa calcula la evolución mediante la regla anterior y después
la compara con esta expresión. La referencia exacta no se utiliza para
rellenar los valores futuros de la malla.

Para representar la oscilación en una gráfica, en cada instante se mide
la amplitud del perfil espacial $\cos(j\widehat x)$ presente en el campo
calculado. Para obtenerla se multiplican los valores del campo por ese
coseno, se suman y se divide por la suma de los cuadrados del coseno.
Esta normalización devuelve amplitud uno para el perfil inicial. En la
solución exacta, la amplitud temporal así obtenida es
$\cos(\widehat\omega_j\widehat t)$. Esta es la magnitud representada en
la figura~\ref{fig:evolucion-plana}.

\begin{figure}[htbp]
\centering
\includegraphics[width=\textwidth]{../results/flat/figures/evolution.pdf}
\caption{Amplitud temporal del modo $j=2$: solución exacta y cálculo con
$N=256$, $\widehat\mu=0.7$ y amplitud inicial uno. El panel inferior
muestra la diferencia entre ambas señales, que apenas se aprecia al
comparar directamente las oscilaciones.}
\label{fig:evolucion-plana}
\end{figure}

Las dos oscilaciones casi coinciden en la escala del panel superior.
La diferencia del panel inferior muestra, sin embargo, un pequeño error
que aumenta en la ventana estudiada: una frecuencia ligeramente distinta
va desplazando los máximos y mínimos de la señal. La comparación visual
se complementa por ello con las mediciones de frecuencia y de error.

\subsection{Frecuencias y relación de dispersión}
La frecuencia se mide a partir de la señal calculada. Se ajusta un coseno
permitiendo variar su frecuencia, amplitud, fase y valor medio, hasta que
reproduzca lo mejor posible los valores registrados. La frecuencia
teórica no se proporciona al ajuste. Los detalles de implementación se
identifican en el apartado~\ref{sec:reproduccion-plana}.

La tabla~\ref{tab:dispersion-plana} compara el resultado con la predicción
$\widehat\omega_j=\sqrt{j^2+0.7^2}$. El error relativo es la diferencia
absoluta entre la frecuencia medida y la exacta, dividida por la exacta.
Por ejemplo, para $j=2$ obtenemos $2.11892519$ frente a $2.11896201$:
un error relativo de aproximadamente $0.00174$ por ciento.

\begin{table}[htbp]
\centering\small
\begin{tabular}{rccc}
\toprule
$j$ & Frecuencia exacta & Frecuencia medida & Error relativo \\
\midrule
1 & 1.22065556 & 1.22066419 & $7.071\times10^{-6}$ \\
2 & 2.11896201 & 2.11892519 & $1.738\times10^{-5}$ \\
3 & 3.08058436 & 3.08039358 & $6.193\times10^{-5}$ \\
4 & 4.06078810 & 4.06028029 & $1.251\times10^{-4}$ \\
\bottomrule
\end{tabular}

\caption{Frecuencias adimensionales medidas para los cuatro perfiles,
con $N=256$ y $0\le\widehat t\le20$. La última columna expresa el error
relativo como fracción, sin multiplicar por cien.}
\label{tab:dispersion-plana}
\end{table}

Los errores relativos de los cuatro casos van de $7.1\times10^{-6}$ a
$1.25\times10^{-4}$. La figura~\ref{fig:dispersion-plana} muestra tanto
el acuerdo con la relación de dispersión como el error que queda oculto
cuando se superponen los valores exactos y los medidos. Para estos
perfiles y esta resolución, el procedimiento recupera la relación entre
frecuencia y número de onda con los errores cuantificados en la tabla.

\begin{figure}[htbp]
\centering
\includegraphics[width=\textwidth]{../results/flat/figures/dispersion.pdf}
\caption{Relación entre frecuencia y número de onda al cuadrado
(izquierda), y error relativo de las frecuencias medidas (derecha).
Todas las magnitudes son adimensionales; $\widehat k_j=j$ y
$\widehat\mu=0.7$.}
\label{fig:dispersion-plana}
\end{figure}

\subsection{Qué ocurre al utilizar más puntos}
Además de medir frecuencias, comparamos el campo en todas las posiciones
de la malla. En cada instante se resta el valor exacto al calculado en
cada punto, se elevan esas diferencias al cuadrado, se suman con peso
$h$ y se toma la raíz cuadrada:
\[
e_n=\left[h\sum_{i=0}^{N-1}
\left(\Phi_i^n-
\widehat\Phi_{\rm ex}(\widehat x_i,\widehat t_n)\right)^2\right]^{1/2}.
\]
Así obtenemos una medida conjunta del error espacial. El peso $h$ hace
que la suma aproxime la integral del error al cuadrado y permite comparar
mallas con distinto número de puntos. Es la norma $L^2$ discreta del
error. La cifra de la tabla~\ref{tab:convergencia-plana} es el mayor valor
de $e_n$ registrado durante cada evolución; no es un porcentaje.

\begin{table}[htbp]
\centering\small
\begin{tabular}{rccc}
\toprule
$N$ & Paso temporal & Error máximo & Orden observado \\
\midrule
64 & 0.078431 & $2.108\times10^{-2}$ & --- \\
128 & 0.039216 & $5.262\times10^{-3}$ & 2.0019 \\
256 & 0.019627 & $1.304\times10^{-3}$ & 2.0122 \\
512 & 0.009814 & $3.261\times10^{-4}$ & 2.0001 \\
\bottomrule
\end{tabular}

\caption{Error del modo $j=2$ al aumentar la resolución. El paso temporal
y el campo están expresados en variables adimensionales. El orden compara
el error de cada fila con el de la anterior.}
\label{tab:convergencia-plana}
\end{table}

Al duplicar el número de puntos, la separación espacial se reduce a la
mitad y el paso temporal también se reduce aproximadamente a la mitad.
Si el error se divide por cuatro, decimos que el orden observado es dos.
Los órdenes de la tabla, entre $2.000$ y $2.012$, son compatibles con ese
comportamiento. Entre $N=64$ y $N=512$, el error máximo disminuye de
aproximadamente $0.02108$ a $0.0003261$.

\begin{figure}[htbp]
\centering
\includegraphics[width=0.88\textwidth]{../results/flat/figures/convergence.pdf}
\caption{Error máximo del campo frente a la separación espacial para
$j=2$. Ambos ejes usan escala logarítmica para comparar las reducciones
del error. La referencia discontinua corresponde a un error proporcional
a $h^2$; se normaliza con la malla más fina y no es otra ejecución.}
\label{fig:convergencia-plana}
\end{figure}

Esta prueba comprueba que el resultado mejora de forma regular con el
refinamiento utilizado. Como se han reducido simultáneamente los pasos
espacial y temporal, no permite atribuir por separado el error a cada uno.

% !TEX root = ../main.tex
\section{Métodos en Schwarzschild}
% T12--T14: justificar por separado contornos espectrales físicos y fronteras
% numéricas del dominio temporal truncado. Ambos métodos están pendientes.

\pendiente{Desarrollar y justificar la evolución en Schwarzschild y el problema radial de autovalores. Sus condiciones físicas y sus fronteras numéricas deben tratarse por separado.}

% !TEX root = ../main.tex
\enlargethispage{\baselineskip}
\section{Alcance de esta primera comprobación}

La primera comprobación respalda el cálculo de los perfiles planos elegidos
y su mejora al refinar conjuntamente espacio y tiempo. Se limita a cosenos
individuales, una masa adimensional y un tiempo final fijo. Otros datos
iniciales, paquetes de ondas, tiempos mayores y fronteras abiertas requieren
pruebas adicionales. Schwarzschild requiere una validación propia.

